Experimental virtual-machine studies combine execution, observation, storage,
and downstream analysis, but these responsibilities are often implemented as
one machine-specific path. This coupling makes it difficult to compare
heterogeneous machines under a shared experimental procedure and to account
for what was observed, delivered, and persisted. GlassVM separates
machine-owned semantics from experimental orchestration. It validates a run
at preparation time, negotiates requested observations before execution, and
delivers typed evidence synchronously to caller-owned sinks. Streaming
normalizers derive selected capabilities from bounded reducer state, while a
reference Rust recorder reports persistence and publication separately from
execution and evidence fulfillment.

We evaluate this design across CHIP-8, PICO-8, and TIC-80, whose execution
models and fidelity boundaries differ. The evaluation covers shared-contract
conformance, selective observation, memory behavior as run length grows,
Python sink delivery, Rust/Python semantic equivalence, package-index-only
consumption, and a small downstream consumer. Final sealed measurements are
not yet available: memory behavior is \TBAMemoryResult, selective observation
cost is \TBAOverheadResult, persistence cost is \TBAStorageResult, and Python
sink overhead is \TBAPythonSinkResult. The manuscript leaves these empirical
claims unresolved until their sealed artifact receipts are available.
